\chapter{Scomporre il monolite}
\label{cap:monolite}

\begin{inbreve}
Molte applicazioni a microservizi nascono da un monolite che, crescendo, è diventato difficile da modificare. Il
capitolo spiega come scomporlo in servizi. Prima si cercano nel codice i \emph{seam}, le linee lungo cui il monolite
si può tagliare, e si decide da dove iniziare. Poi si affronta la parte più difficile, il database: si spezzano le
dipendenze tra tabelle, si gestiscono le transazioni che coinvolgono più servizi e si ricostruisce il reporting.
Presuppone il capitolo~\ref{cap:microservizi}.
\end{inbreve}

\section{Il monolite cresce}
\label{sec:split-monolite}

Si richiamano i microservizi (capitolo~\ref{cap:microservizi}). Si sviluppano le applicazioni come insiemi di servizi
stateless di piccole dimensioni per due motivi (\rif{sec:micro-motivi}):
\begin{enumerate}
  \item ridurre il \emph{lead time} di nuove feature e aggiornamenti, cioè il tempo tra la decisione di fare una
    modifica e la sua distribuzione agli utenti;
  \item scalare in modo efficace.
\end{enumerate}
I servizi girano ciascuno nel proprio container, comunicano con meccanismi leggeri e sono costruiti attorno a
capacità di business. Decentralizzano la gestione dei dati, sono distribuibili in modo indipendente, scalabili
orizzontalmente e resilienti ai guasti. Il tutto con un approccio DevOps (\rif{sec:micro-distribuzione}).

Un monolite invece cresce nel tempo: vi si aggiungono di continuo nuove funzionalità e nuove righe di codice. Con la
crescita perde due proprietà (\rif{sec:micro-microservizi}):
\begin{itemize}
  \item \textbf{la cohesion}: codice di ogni genere, senza relazione, finisce nello stesso blocco;
  \item \textbf{il loose coupling}: un piccolo cambiamento può avere effetti sul resto del monolite, e per
    distribuirlo bisogna ridistribuire l'intero blocco.
\end{itemize}

\begin{idea}[Quando scomporre]
\uline{Il monolite si scompone \emph{quando} diventa un problema}, non prima. Far crescere un sistema finché va diviso
è normale. Bisogna però accorgersi del momento giusto e dividerlo prima che la divisione costi troppo
(\rif{sec:split-conclusioni}).
\end{idea}

\section{Trovare i seam}
\label{sec:split-seam}

Il monolite si scompone lungo i suoi \emph{seam}. In sartoria il \ingl{seam} è la cucitura, la linea lungo cui due pezzi
di tessuto sono cuciti insieme in un abito.

\begin{definizione}[Seam]
Un \termine{seam} è una porzione di codice che si può trattare in isolamento e su cui si può lavorare senza effetti sul
resto del codice.
\end{definizione}

L'obiettivo è trovare i seam che possono diventare \emph{confini di servizio}. Due strumenti aiutano.
\begin{itemize}
  \item La nozione di \termine{bounded context}. Alcune informazioni di un contesto non devono essere comunicate
    all'esterno, altre si possono condividere con altri contesti. È come una cellula: la membrana stabilisce che cosa
    sta dentro, che cosa sta fuori e che cosa può passare.
  \item I concetti di \emph{namespace} dei linguaggi di programmazione, per esempio i \emph{package} di Java.
\end{itemize}

\subsubsection{Che cosa sta dentro e che cosa passa la membrana\aggiunta}
Un bounded context è una parte dell'organizzazione con una responsabilità precisa, per esempio il magazzino. Al suo
interno usa molte informazioni che interessano solo a lui: in quale scaffale sta un articolo, quali ordini sono in
preparazione. Queste restano dentro. Altre informazioni servono anche fuori: la contabilità, per esempio, deve sapere
quanti articoli sono stati spediti. Queste passano la membrana, attraverso un'interfaccia esplicita. Un contesto così
delimitato è un buon candidato a diventare un servizio: l'interfaccia del contesto diventa l'API del servizio.

\begin{esempio}[MusicCorp]
MusicCorp è un rivenditore online di musica. Il suo monolite copre quattro contesti:
\begin{description}[leftmargin=0pt, style=unboxed]
  \item[Catalog.] Tutto ciò che riguarda i metadati degli articoli in vendita.
  \item[Finance.] Contabilità e rendiconti, pagamenti, rimborsi, ecc.
  \item[Warehouse.] Spedizione e reso degli ordini dei clienti, gestione dei livelli di magazzino, ecc.
  \item[Recommendation.] Un sistema di raccomandazione rivoluzionario, in attesa di brevetto.
\end{description}
\end{esempio}

\subsection{Organizzare il codice attorno ai seam}
\label{sec:split-package}

Individuati i contesti, si \textbf{creano dei package che li rappresentano e poi vi si sposta il codice esistente}.
\begin{itemize}
  \item Il refactoring e i test procedono in modo incrementale.
  \item Il codice che resta fuori da tutti i package può rivelare un nuovo bounded context.
  \item Il codice stesso serve ad analizzare le dipendenze tra i package. I package che rappresentano i bounded context
    dovrebbero interagire come i gruppi reali dell'organizzazione. Per esempio, il package \emph{Warehouse} non
    dovrebbe dipendere dal package \emph{Finance} se nell'organizzazione reale non c'è una dipendenza del genere.
  \item Strumenti come \emph{Structure 101} mostrano graficamente le dipendenze tra package. Nel grafo ogni nodo è un
    package e ogni arco una dipendenza, con il numero di riferimenti da un package all'altro.
\end{itemize}
\uline{Tutto questo si fa in modo incrementale, davvero}: un passo alla volta, verificando ogni passo con i test.

\section{Da dove iniziare}
\label{sec:split-inizio}

Si inizia dal seam che porta il \textbf{massimo beneficio}. I criteri tipici (\ingl{drivers}) sono cinque; gli esempi
sono quelli di MusicCorp.
\begin{description}[leftmargin=0pt, style=unboxed]
  \item[Pace of change.] Il ritmo dei cambiamenti. Se sta per arrivare una mole di modifiche alla gestione del
    magazzino, conviene separare per primo il seam \emph{Warehouse}. Così le modifiche si faranno su un servizio
    piccolo e distribuibile da solo.
  \item[Team structure.] La struttura dei team. Se il team di sviluppo è diviso tra due regioni geografiche, conviene
    separare le parti su cui lavora ciascuna delle due metà.
  \item[Security.] Se tutte le informazioni sensibili sono gestite dal codice di \emph{Finance}, separarlo permette di
    proteggere di più solo quel servizio.
  \item[Technology.] Se il team del sistema di \emph{Recommendation} sperimenta nuovi algoritmi in Clojure, un servizio
    separato può usare un linguaggio diverso dal resto.
  \item[Tangled dependencies.] Le dipendenze aggrovigliate. Se possibile, si estrae per primo il seam da cui gli altri
    dipendono meno, guardando per esempio il grafo delle dipendenze.
\end{description}

\begin{esempio}[Il seam meno dipeso\aggiunta]
Si supponga che il grafo delle dipendenze tra i package di MusicCorp sia quello di \figrif{fig:split-dipendenze}. Le
dipendenze sono ipotetiche. Da Catalog dipendono tre package: estrarlo per primo costringe a modificarli tutti e tre.
Da Recommendation e da Finance non dipende nessuno. Recommendation però dipende da un solo package, Finance da due.
Quindi Recommendation è il seam più facile da estrarre: basta trasformare in chiamate all'API di Catalog le sue
chiamate a Catalog.
\end{esempio}

\begin{figure}[ht]
\centering
\begin{tikzpicture}
\node[blocco, minimum width=2.6cm] (cat) at (0,0) {Catalog};
\node[blocco, minimum width=2.6cm] (fin) at (-4.2,1.6) {Finance};
\node[blocco, minimum width=2.6cm] (war) at (-4.2,-1.6) {Warehouse};
\node[blocco, minimum width=2.6cm] (rec) at (4.2,0) {Recommendation};
\draw[freccia] (fin) -- (cat);
\draw[freccia] (war) -- (cat);
\draw[freccia] (rec) -- (cat);
\draw[freccia] (fin) -- (war);
\node[font=\footnotesize, anchor=north] at (0,-0.5) {usato da 3 package};
\node[font=\footnotesize, anchor=north] at (4.2,-0.5) {usato da nessuno};
\node[font=\footnotesize, anchor=south] at (-4.2,2.05) {usato da nessuno};
\node[font=\footnotesize, anchor=north] at (-4.2,-2.05) {usato da 1 package};
\end{tikzpicture}
\caption{Grafo ipotetico delle dipendenze tra i package di MusicCorp. Una freccia da A a B significa che A usa B.
Recommendation non è usato da nessuno e usa un solo package: è il primo candidato a diventare un servizio.}
\label{fig:split-dipendenze}
\end{figure}

\section{Il database}
\label{sec:split-database}

Il database è la madre di tutte le dipendenze aggrovigliate. Spesso è usato per \emph{integrare} più servizi: invece
di chiamarsi a vicenda, le parti del sistema leggono e scrivono le stesse tabelle. Bisogna quindi trovare i seam
anche nel database, per poterlo dividere in modo pulito. Tre attività aiutano.
\begin{itemize}
  \item \textbf{Capire quali parti del codice leggono e scrivono il database}, per esempio con strumenti come
    Hibernate. Hibernate è una libreria Java di \emph{object-relational mapping}: associa ogni classe alla tabella in
    cui sono salvati i suoi oggetti. Le sue associazioni mostrano quale codice usa quali tabelle.
  \item \textbf{Individuare i vincoli a livello di database}, per esempio relazioni di foreign key tra tabelle usate
    da parti diverse del codice. Strumenti come SchemaSpy disegnano lo schema con le tabelle e le relazioni.
  \item \textbf{Notare le tabelle usate da più bounded context.}
\end{itemize}

Una \termine{foreign key} è una colonna di una tabella che contiene la chiave di una riga di un'altra tabella. Il
database garantisce che il valore esista davvero nell'altra tabella: non si può inserire un riferimento a una riga
inesistente, né cancellare una riga ancora riferita.

\section{Refactoring del database}
\label{sec:split-refactoring}

Quattro situazioni tipiche, tutte nel monolite di MusicCorp.

\subsection{Spezzare una relazione di foreign key}
\label{sec:split-fk}

Il codice di \emph{Catalog} usa la tabella \emph{Line items} per memorizzare le informazioni su un album. Il codice di
\emph{Finance} usa la tabella \emph{Ledger}, il libro contabile, per tracciare le informazioni finanziarie. Ogni mese
Finance genera report come ``abbiamo venduto 600 copie di \emph{Dance Monkey} di Tones and I e incassato 2\,000
USD''. Per scriverli:
\begin{itemize}
  \item il codice di Finance legge la tabella Line items, che appartiene a Catalog;
  \item esiste una relazione di foreign key dalla tabella Ledger alla tabella Line items.
\end{itemize}
Il problema è: \emph{come si impedisce al codice di Finance di accedere alla tabella Line items?}

\textbf{Soluzione}: il package Catalog espone i dati tramite un'API (\figrif{fig:split-fk}). Finance chiede a Catalog i
dati dell'album invece di leggere la sua tabella. Le conseguenze sono due.
\begin{itemize}
  \item Si introduce un \textbf{overhead}: una chiamata all'API costa più di una lettura dal database.
  \item Si \textbf{perde la relazione di foreign key}. I vincoli vanno gestiti nei servizi invece che a livello di
    database. Può servire un controllo di coerenza tra servizi, oppure azioni che ripuliscano i dati collegati.
\end{itemize}

\begin{attenzione}[Una scelta che non è tecnica]
Se e come controllare la coerenza spesso non è una scelta tecnologica. Per esempio: come si gestisce un ordine che
si riferisce all'ID di catalogo di un articolo rimosso dal catalogo? Lo decide il business, non il database.
\end{attenzione}

\begin{figure}[ht]
\centering
\begin{tikzpicture}
% prima
\node[font=\small\bfseries\sffamily] at (1.6,2.75) {MusicCorpMono};
\draw (-0.8,1.0) rectangle (4.0,3.1);
\node[sv, minimum width=1.8cm] (c) at (0.3,1.7) {Catalog};
\node[sv, minimum width=1.8cm] (f) at (2.9,1.7) {Finance};
\tabdb{0.3}{-0.3}{Line items}{li}
\tabdb{2.9}{-0.3}{Ledger}{le}
\draw[->] (c) -- (li);
\draw[->] (f) -- (le);
\draw[->, dashed] (f.south west) -- (li.north east);
\draw[->, dashed] (le.west) -- (li.east);
% freccia
\draw[->, very thick, verde] (4.5,0.8) -- (5.6,0.8);
% dopo
\node[sv, minimum width=2.4cm] (c2) at (6.9,1.7) {Catalog service};
\node[sv, minimum width=2.4cm] (f2) at (10.7,1.7) {Finance service};
\tabdb{6.9}{-0.3}{Line items}{li2}
\tabdb{10.7}{-0.3}{Ledger}{le2}
\draw[->] (c2) -- (li2);
\draw[->] (f2) -- (le2);
\draw[->, dashed] (f2) -- node[above, font=\footnotesize] {API} (c2);
\end{tikzpicture}
\caption{Spezzare una foreign key. A sinistra, nel monolite, Finance legge la tabella Line items di Catalog (freccia
tratteggiata in alto) e Ledger ha una foreign key verso Line items (freccia tratteggiata in basso). A destra Finance
chiede i dati all'API del servizio Catalog: ogni servizio usa solo le proprie tabelle.}
\label{fig:split-fk}
\end{figure}

\subsection{Dati statici condivisi}
\label{sec:split-statici}

Catalog, Warehouse e Finance leggono la stessa tabella \emph{Country codes}, che contiene i codici dei paesi. Sono
dati \emph{statici}: cambiano raramente. Le soluzioni sono tre (\figrif{fig:split-statici}).
\begin{enumerate}
  \item \textbf{Duplicare la tabella per ogni package.} Si pone un possibile problema di coerenza: un aggiornamento va
    fatto in tutte le copie.
  \item \textbf{Trattare i dati statici condivisi come codice}, per esempio come un file di configurazione. I problemi
    di coerenza restano, anche se forse è più facile aggiornare un file di configurazione che una tabella del
    database. È la soluzione preferita.
  \item \textbf{Spostare i dati statici in un servizio separato.} Per dati così semplici rischia di essere eccessivo
    (\ingl{overkill}).
\end{enumerate}

\begin{figure}[ht]
\centering
\begin{tikzpicture}
\node[font=\small\bfseries\sffamily] at (2.4,2.75) {MusicCorpMono};
\draw (-0.6,1.0) rectangle (5.4,3.1);
\node[sv, minimum width=1.6cm] (c) at (0.4,1.7) {Catalog};
\node[sv, minimum width=1.6cm] (w) at (2.4,1.7) {Warehouse};
\node[sv, minimum width=1.6cm] (f) at (4.4,1.7) {Finance};
\tabdb{2.4}{-0.3}{Country codes}{cc}
\draw[->] (c.south) -- (cc.north west);
\draw[->] (w.south) -- (cc.north);
\draw[->] (f.south) -- (cc.north east);
% soluzioni
\node[font=\small, anchor=west] at (6.6,2.7) {(1) una copia per package};
\foreach \x/\n in {6.9/C, 8.1/W, 9.3/F} {
  \node[sv, minimum width=1.0cm, minimum height=0.6cm] (s\n) at (\x,2.0) {\n};
  \draw (\x-0.4,1.05) rectangle (\x+0.4,1.4);
  \draw[->] (s\n) -- (\x,1.4);}
\node[font=\small, anchor=west] at (6.6,0.6) {(2) dati come configurazione};
\node[draw, font=\footnotesize, anchor=west] at (6.7,0.0) {countries.conf};
\node[font=\footnotesize, anchor=west] at (9.0,0.0) {nel codice};
\node[font=\small, anchor=west] at (6.6,-0.8) {(3) servizio separato};
\node[sv, minimum width=2.6cm, anchor=west] at (6.7,-1.4) {Country service};
\end{tikzpicture}
\caption{Dati statici condivisi. A sinistra tre package leggono la stessa tabella dei codici dei paesi. A destra le tre
soluzioni: una copia della tabella per ogni package (C, W, F), un file di configurazione distribuito con il codice, un
servizio dedicato.}
\label{fig:split-statici}
\end{figure}

\subsection{Dati modificabili condivisi}
\label{sec:split-mutabili}

Il codice di \emph{Finance} registra i pagamenti e i rimborsi dei clienti. Il codice di \emph{Warehouse} aggiorna gli
ordini spediti e ricevuti. Entrambi leggono e \emph{scrivono} la tabella \emph{Customer record}. Il concetto di
dominio \emph{Customer} è modellato solo in modo implicito, dentro il database.

\textbf{Soluzione}: si crea un nuovo servizio \emph{Customer} e si fa in modo che Finance e Warehouse invochino la sua
API (\figrif{fig:split-mutabili}). Il concetto implicito diventa un servizio esplicito, unico responsabile dei dati
dei clienti.

\begin{figure}[ht]
\centering
\begin{tikzpicture}
\node[font=\small\bfseries\sffamily] at (1.3,2.75) {MusicCorp};
\draw (-0.5,1.0) rectangle (3.1,3.1);
\node[sv, minimum width=1.5cm] (w) at (0.3,1.7) {Warehouse};
\node[sv, minimum width=1.5cm] (f) at (2.3,1.7) {Finance};
\tabdb{1.3}{-0.3}{Customer record}{cr}
\draw[->] (w.south) -- (cr.north west);
\draw[->] (f.south) -- (cr.north east);
\draw[->, very thick, verde] (3.8,0.8) -- (4.9,0.8);
\node[font=\small\bfseries\sffamily] at (8.0,2.75) {MusicCorp system};
\node[sv, minimum width=2.4cm] (f2) at (6.6,1.9) {Finance service};
\node[sv, minimum width=2.4cm] (w2) at (9.4,1.9) {Warehouse service};
\node[sv, minimum width=2.4cm] (cu) at (8.0,0.0) {Customer service};
\draw[->] (f2) -- (cu);
\draw[->] (w2) -- (cu);
\end{tikzpicture}
\caption{Dati modificabili condivisi. A sinistra Warehouse e Finance scrivono la stessa tabella dei clienti. A destra il
nuovo servizio Customer gestisce i dati dei clienti e gli altri due servizi lo invocano.}
\label{fig:split-mutabili}
\end{figure}

\subsection{Tabelle condivise}
\label{sec:split-tabelle}

\emph{Catalog} deve memorizzare il nome e il prezzo dei dischi in vendita. \emph{Warehouse} deve tenere la registrazione
elettronica dell'inventario. I due concetti sono memorizzati nella stessa tabella \emph{Item}.

\textbf{Soluzione}: si memorizzano i due concetti separatamente, dividendo la tabella in due: \emph{Catalog item} per
Catalog e \emph{Stock levels} (livelli di magazzino) per Warehouse (\figrif{fig:split-tabelle}).

\begin{attenzione}[Due casi da non confondere\aggiunta]
Nei dati modificabili condivisi (\rif{sec:split-mutabili}) due package usano \emph{lo stesso concetto}, il cliente: si
crea un servizio per quel concetto. Nelle tabelle condivise due package usano \emph{concetti diversi} finiti nella
stessa tabella: si divide la tabella, un concetto per parte.
\end{attenzione}

\begin{figure}[ht]
\centering
\begin{tikzpicture}
\node[font=\small\bfseries\sffamily] at (1.3,2.75) {MusicCorp};
\draw (-0.5,1.0) rectangle (3.1,3.1);
\node[sv, minimum width=1.5cm] (c) at (0.3,1.7) {Catalog};
\node[sv, minimum width=1.5cm] (w) at (2.3,1.7) {Warehouse};
\tabdb{1.3}{-0.3}{Item}{it}
\draw[->] (c.south) -- (it.north west);
\draw[->] (w.south) -- (it.north east);
\draw[->, very thick, verde] (3.8,0.8) -- (4.9,0.8);
\node[font=\small\bfseries\sffamily] at (7.7,2.75) {MusicCorp};
\draw (5.6,1.0) rectangle (9.8,3.1);
\node[sv, minimum width=1.5cm] (c2) at (6.6,1.7) {Catalog};
\node[sv, minimum width=1.5cm] (w2) at (8.8,1.7) {Warehouse};
\tabdb{6.6}{-0.3}{Catalog item}{ci}
\tabdb{8.8}{-0.3}{Stock levels}{sl}
\draw[->] (c2) -- (ci);
\draw[->] (w2) -- (sl);
\end{tikzpicture}
\caption{Tabelle condivise. A sinistra nome e prezzo (Catalog) e inventario (Warehouse) stanno nella stessa tabella Item.
A destra la tabella è divisa: ogni package usa la tabella del proprio concetto.}
\label{fig:split-tabelle}
\end{figure}

\subsection{Prima gli schemi o prima i servizi?}
\label{sec:split-ordine}

Il passaggio dal monolite ai servizi si può fare in tre tappe (\figrif{fig:split-ordine}): uno schema unico; schemi
separati con il codice ancora nel monolite; applicazione divisa in servizi.
\begin{itemize}
  \item Separare gli schemi può aumentare il numero di chiamate al database. Può anche rompere l'integrità
    transazionale (\rif{sec:split-transazioni}).
  \item \textbf{Separare gli schemi tenendo insieme il codice applicativo rende più facile tornare indietro.} Se la
    separazione degli schemi crea problemi, si annulla senza aver toccato la struttura del codice.
\end{itemize}

\begin{figure}[ht]
\centering
\begin{tikzpicture}
% 1
\node[font=\footnotesize] at (1.9,3.15) {Monolithic service};
\draw (0.1,1.55) rectangle (3.7,3.45);
\node[sv, minimum width=1.7cm, align=center] (a1) at (0.95,2.3) {Finance\\code};
\node[sv, minimum width=1.7cm, align=center] (a2) at (2.85,2.3) {Customer\\code};
\node[dbs, minimum width=2.0cm] (ad) at (1.9,0.1) {Monolithic\\schema};
\draw (a1.south) |- (1.9,1.2) -| (a2.south);
\draw[->] (1.9,1.2) -- (ad.north);
\node[font=\footnotesize] at (1.9,-1.05) {1. Single schema};
\draw[grigio] (4.25,-1.4) -- (4.25,3.55);
% 2
\node[font=\footnotesize] at (6.6,3.15) {Monolithic service};
\draw (4.8,1.55) rectangle (8.4,3.45);
\node[sv, minimum width=1.7cm, align=center] (b1) at (5.65,2.3) {Finance\\code};
\node[sv, minimum width=1.7cm, align=center] (b2) at (7.55,2.3) {Customer\\code};
\node[dbs, minimum width=1.7cm] (bd1) at (5.65,0.1) {Finance\\schema};
\node[dbs, minimum width=1.7cm] (bd2) at (7.55,0.1) {Customer\\schema};
\draw[->] (b1) -- (bd1); \draw[->] (b2) -- (bd2);
\node[font=\footnotesize] at (6.6,-1.05) {2. Split schemas};
\draw[grigio] (8.95,-1.4) -- (8.95,3.55);
% 3
\node[sv, minimum width=1.7cm, align=center] (c1) at (10.35,2.3) {Finance\\service};
\node[sv, minimum width=1.7cm, align=center] (c2) at (12.25,2.3) {Customer\\service};
\node[dbs, minimum width=1.7cm] (cd1) at (10.35,0.1) {Finance\\schema};
\node[dbs, minimum width=1.7cm] (cd2) at (12.25,0.1) {Customer\\schema};
\draw[->] (c1) -- (cd1); \draw[->] (c2) -- (cd2);
\node[font=\footnotesize, align=center] at (11.3,-1.15) {3. Split application\\into services};
\end{tikzpicture}
\caption{Le tre tappe della scomposizione. 1: il servizio monolitico usa uno schema unico. 2: gli schemi sono separati, il
codice è ancora nel monolite. 3: anche il codice è diviso in servizi, ciascuno con il proprio schema.}
\label{fig:split-ordine}
\end{figure}

\section{Transazioni}
\label{sec:split-transazioni}

Le transazioni \emph{tutto-o-niente} (le transazioni ACID, \rif{sec:micro-dati}) sono utili: aggiornano più tabelle
``in un colpo solo'' e lasciano il database sempre in uno stato coerente.

\begin{esempio}[Creare un nuovo ordine]
Quando si crea un nuovo ordine si aggiornano due tabelle:
\begin{itemize}
  \item la tabella \emph{Order}, con l'ordine (codice di Customer);
  \item la tabella \emph{Picking}, con l'articolo da prelevare per la spedizione (codice di Warehouse).
\end{itemize}
Con uno schema monolitico le due insert fanno parte di un'unica transazione. Con due schemi separati ogni insert ha il
proprio confine transazionale, e la sicurezza transazionale si perde (\figrif{fig:split-transazioni}).
\end{esempio}

La domanda è: \emph{che cosa succede se l'inserimento dell'ordine riesce e quello del prelievo fallisce?} Il cliente ha
un ordine registrato, ma nessuno spedirà l'articolo. Le soluzioni sono tre.

\begin{figure}[ht]
\centering
\begin{tikzpicture}
% monolitico: servizio in alto, schema unico in basso, un solo confine
\node[font=\footnotesize] at (2.4,4.15) {Monolithic service};
\draw (0.2,2.55) rectangle (4.6,4.45);
\node[sv, minimum width=1.9cm, align=center] (a1) at (1.2,3.3) {Customer\\code};
\node[sv, minimum width=1.9cm, align=center] (a2) at (3.6,3.3) {Warehouse\\code};
\draw (0.2,-0.95) rectangle (4.6,1.0);
\node[blocco, minimum width=1.5cm, align=center, font=\footnotesize] (t1) at (1.2,0.35) {Order\\table};
\node[blocco, minimum width=1.5cm, align=center, font=\footnotesize] (t2) at (3.6,0.35) {Picking\\table};
\node[font=\footnotesize] at (2.4,-0.5) {Monolithic schema};
\draw[->] (a1) -- node[right, pos=0.3, font=\footnotesize, align=left] {insert\\order} (t1);
\draw[->] (a2) -- node[left, pos=0.3, font=\footnotesize, align=right] {insert\\picking} (t2);
\draw[rosso, dashed, thick] (0.0,-1.15) rectangle (4.8,1.55);
\node[font=\footnotesize, text=rosso] at (2.4,-1.5) {un solo confine transazionale};
% schemi separati: due confini
\node[font=\footnotesize] at (8.8,4.15) {Monolithic service};
\draw (6.4,2.55) rectangle (11.2,4.45);
\node[sv, minimum width=1.9cm, align=center] (b1) at (7.5,3.3) {Customer\\code};
\node[sv, minimum width=1.9cm, align=center] (b2) at (10.1,3.3) {Warehouse\\code};
\draw (6.4,-0.95) rectangle (8.6,1.0);
\draw (9.0,-0.95) rectangle (11.2,1.0);
\node[blocco, minimum width=1.5cm, align=center, font=\footnotesize] (u1) at (7.5,0.35) {Order\\table};
\node[blocco, minimum width=1.5cm, align=center, font=\footnotesize] (u2) at (10.1,0.35) {Picking\\table};
\node[font=\footnotesize, align=center] at (7.5,-0.5) {Customer\\schema};
\node[font=\footnotesize, align=center] at (10.1,-0.5) {Warehouse\\schema};
\draw[->] (b1) -- node[right, pos=0.3, font=\footnotesize, align=left] {insert\\order} (u1);
\draw[->] (b2) -- node[left, pos=0.3, font=\footnotesize, align=right] {insert\\picking} (u2);
\draw[rosso, dashed, thick] (6.25,-1.15) rectangle (8.75,1.55);
\draw[rosso, dashed, thick] (8.85,-1.15) rectangle (11.35,1.55);
\node[font=\footnotesize, text=rosso] at (8.8,-1.5) {due confini transazionali separati};
\end{tikzpicture}
\caption{Transazioni e schemi. A sinistra le due insert stanno in un unico confine transazionale (tratteggio rosso):
riescono entrambe o falliscono entrambe. A destra, con due schemi, ogni insert ha il proprio confine: una può
riuscire mentre l'altra fallisce.}
\label{fig:split-transazioni}
\end{figure}

\subsection{Soluzione 1: transazioni compensative}

Si \textbf{annulla l'intera operazione con una o più transazioni compensative}, come nel pattern Saga
(\rif{sec:micro-dati}). Per esempio, se l'inserimento dell'ordine riesce e quello del prelievo fallisce:
\begin{itemize}
  \item si avvia una transazione compensativa che annulla la transazione già confermata nella tabella Order;
  \item si segnala il fallimento attraverso l'interfaccia utente.
\end{itemize}
Restano dei problemi.
\begin{itemize}
  \item Dove sta la logica che gestisce le transazioni compensative? Nel servizio Customer, nel servizio Warehouse o
    altrove?
  \item Che cosa succede se fallisce la transazione compensativa? Si riprova (\ingl{retry}).
  \item Che cosa succede se bisogna garantire la coerenza di 5 operazioni invece di 2? Ogni operazione può fallire
    dopo che le precedenti sono riuscite, e per ciascun caso serve la giusta sequenza di compensazioni.
\end{itemize}

\subsection{Soluzione 2: transazioni distribuite}

Si usano le \termine{transazioni distribuite}. Un \termine{transaction manager} orchestra le varie transazioni locali.
L'algoritmo più comune è il \emph{two-phase commit}.

\begin{algoritmo}[Two-phase commit]
\begin{enumerate}
  \item \textbf{Fase di voto} (\ingl{voting phase}). Ogni partecipante dice al transaction manager se la propria
    transazione locale può andare avanti.
  \item \textbf{Fase di commit.} Se il transaction manager riceve un voto \emph{sì} da tutti i partecipanti, dice a
    tutti di procedere e fare il commit. Altrimenti invia un \emph{rollback} a tutti.
\end{enumerate}
\end{algoritmo}

Il two-phase commit ha delle vulnerabilità.
\begin{itemize}
  \item Se il transaction manager si guasta, le transazioni in sospeso non si completano mai.
  \item Se un partecipante non risponde durante il voto, tutto si blocca.
  \item Che cosa succede se un commit fallisce dopo il voto?
\end{itemize}

\begin{esempio}[Il nuovo ordine con il two-phase commit\aggiunta]
I partecipanti sono lo schema Customer e lo schema Warehouse (\figrif{fig:split-2pc}). Il transaction manager chiede
a entrambi se possono inserire la propria riga. Ciascuno prepara l'inserimento senza renderlo definitivo, blocca le
risorse coinvolte e risponde \emph{sì}. Ricevuti due \emph{sì}, il manager invia a entrambi il \emph{commit}. Se lo
schema Warehouse avesse risposto \emph{no}, il manager avrebbe inviato a entrambi il \emph{rollback}: l'ordine non
sarebbe stato registrato. Tra il voto e il commit le risorse restano bloccate. Per questo un manager guasto o un
partecipante che non risponde bloccano tutti.
\end{esempio}

\begin{figure}[ht]
\centering
\begin{tikzpicture}
\node[blocco, minimum width=3.0cm] (m) at (0,0) {Transaction manager};
\node[blocco, minimum width=3.0cm] (p1) at (5,0) {Customer schema};
\node[blocco, minimum width=3.0cm] (p2) at (10,0) {Warehouse schema};
\foreach \x in {0,5,10} \draw[grigio] (\x,-0.45) -- (\x,-6.3);
\node[font=\footnotesize\bfseries, anchor=east] at (-0.2,-1.3) {Fase 1: voto};
\draw[->] (0,-1.1) -- node[above, font=\footnotesize] {puoi fare il commit?} (5,-1.1);
\draw[->] (0,-1.9) -- node[above, font=\footnotesize, pos=0.75] {puoi fare il commit?} (10,-1.9);
\draw[->, dashed] (5,-2.5) -- node[above, font=\footnotesize] {sì} (0,-2.5);
\draw[->, dashed] (10,-3.1) -- node[above, font=\footnotesize, pos=0.25] {sì} (0,-3.1);
\draw[grigio, densely dotted] (-2.4,-3.55) -- (10.6,-3.55);
\node[font=\footnotesize\bfseries, anchor=east] at (-0.2,-4.2) {Fase 2: commit};
\draw[->] (0,-4.0) -- node[above, font=\footnotesize] {commit} (5,-4.0);
\draw[->] (0,-4.6) -- node[above, font=\footnotesize, pos=0.75] {commit} (10,-4.6);
\draw[->, dashed] (5,-5.3) -- node[above, font=\footnotesize] {fatto} (0,-5.3);
\draw[->, dashed] (10,-5.9) -- node[above, font=\footnotesize, pos=0.25] {fatto} (0,-5.9);
\end{tikzpicture}
\caption{Two-phase commit per il nuovo ordine. Il tempo scorre verso il basso. Nella fase 1 il transaction manager raccoglie
i voti dei due partecipanti; avendo ricevuto due \emph{sì}, nella fase 2 ordina a entrambi il commit. Con anche un solo
\emph{no} avrebbe ordinato a entrambi il rollback.}
\label{fig:split-2pc}
\end{figure}

\subsection{Soluzione 3: riprovare più tardi}

Si \textbf{riprova più tardi} (\ingl{try again later}) e si accetta la \emph{eventual consistency}
(\rif{sec:micro-dati}). È la soluzione preferita.

\subsubsection{Come funziona\aggiunta}
L'ordine viene registrato. L'inserimento nella tabella Picking che è fallito non annulla l'ordine: viene messo in una
coda e ritentato più tardi, finché riesce. Per un intervallo di tempo il sistema è incoerente: l'ordine c'è, il
prelievo non ancora. Quando il tentativo riesce, i dati tornano coerenti. La soluzione funziona se riprovare può
davvero risolvere il problema, come nel caso di un servizio momentaneamente non disponibile.

\section{Reporting}
\label{sec:split-reporting}

Il \termine{reporting} di solito deve mettere insieme dati di molte parti dell'organizzazione per produrre un risultato
utile. Per esempio:
\begin{itemize}
  \item combinare i dati del libro contabile con le descrizioni degli articoli venduti, prese dal catalogo;
  \item combinare gli storici degli acquisti con i profili dei clienti, per capire il comportamento dei clienti.
\end{itemize}

Nei monoliti i report girano spesso su una \emph{read replica} del database primario, per ragioni di prestazioni
(\figrif{fig:split-replica}). La read replica è una copia del database usata solo in lettura e sincronizzata
periodicamente con il database primario. Così le interrogazioni pesanti dei report non rallentano il servizio. Ma
\emph{che cosa si fa se il database è stato diviso?} Le soluzioni sono tre.

\begin{figure}[ht]
\centering
\begin{tikzpicture}
\node[sv, align=center] (ms) at (0,2.2) {Monolithic\\service};
\node[sv, align=center] (rs) at (5,2.2) {Reporting\\system};
\node[dbs] (pd) at (0,0) {Primary\\database};
\node[dbs] (rr) at (5,0) {Reporting\\replica};
\draw[->] (ms) -- node[left, font=\footnotesize] {read/write} (pd);
\draw[->] (rs) -- node[right, font=\footnotesize] {read} (rr);
\draw[->] (pd) -- node[below, font=\footnotesize, align=center] {sincronizzazione\\periodica} (rr);
\end{tikzpicture}
\caption{Reporting in un monolite. Il servizio legge e scrive il database primario. Il sistema di reporting legge una
replica, sincronizzata periodicamente con il primario.}
\label{fig:split-replica}
\end{figure}

\subsection{Soluzione 1: recupero dei dati con chiamate ai servizi}

Il sistema di reporting \emph{preleva} (\ingl{pull}) i dati dai sistemi di origine con chiamate alle loro API.
\begin{itemize}
  \item Non scala bene con grandi volumi di dati.
  \item Le API esposte possono non essere pensate per il reporting: per esempio può non esserci modo di ottenere tutti
    i clienti. Recuperarli uno per uno è inefficiente e può generare carico sul servizio interrogato.
\end{itemize}

\subsection{Soluzione 2: data pump}

Un \termine{data pump} \emph{spinge} (\ingl{push}) i dati verso il sistema di reporting (\figrif{fig:split-pump}).
\begin{itemize}
  \item Il data pump trasforma i dati del database del servizio nello schema del database di reporting.
  \item Il data pump conosce lo schema interno del database del servizio: è un \emph{coupling}. È un prezzo che vale
    la pena pagare per rendere più facile il reporting. Per tenerlo sotto controllo:
    \begin{itemize}
      \item il data pump è costruito e gestito dal team che gestisce il servizio;
      \item il data pump è sotto controllo di versione insieme al servizio.
    \end{itemize}
  \item Il data pump si può appoggiare alle operazioni di backup: è il caso dei \emph{backup data pumps} di Netflix.
    Il backup che il servizio produce comunque diventa la sorgente dei dati per il reporting.
\end{itemize}

\begin{figure}[ht]
\centering
\begin{tikzpicture}
\node[sv, align=center, minimum width=2.8cm] (cs) at (0,2.2) {Customer service\\customer-code v123};
\node[dbs] (db) at (0,0) {Database};
\node[sv, align=center, minimum width=2.8cm] (dp) at (4.0,0) {Customer data pump\\customer-code v123};
\node[sv, align=center] (rt) at (8.0,2.2) {Reporting\\tools};
\node[dbs] (rd) at (8.0,0) {Central\\reporting\\database};
\draw[->] (cs) -- node[left, font=\footnotesize] {read/write} (db);
\draw (db) -- (dp);
\draw[->] (dp) -- (rd);
\draw[->] (rt) -- node[right, font=\footnotesize] {SQL o simili} (rd);
\draw[blu, dashed, thick, rounded corners=8pt] (-1.8,-1.0) rectangle (5.8,3.0);
\node[font=\footnotesize, text=blu, anchor=north east] at (5.7,2.95) {gestiti dallo stesso team};
\end{tikzpicture}
\caption{Data pump. Il data pump legge il database del servizio Customer e scrive nel database centrale di reporting, su cui
lavorano gli strumenti di reporting. Servizio e data pump sono versionati insieme (stessa versione, v123) e li
gestisce lo stesso team (riquadro tratteggiato).}
\label{fig:split-pump}
\end{figure}

\subsection{Soluzione 3: event data pump}

I microservizi possono emettere \emph{eventi} quando cambia lo stato delle entità che gestiscono. Per esempio il servizio
Customer può emettere un evento quando un cliente è creato, aggiornato o cancellato. Si scrive allora un
\emph{subscriber} degli eventi che riversa i dati nel database di reporting: è l'\termine{event data pump}
(\figrif{fig:split-event}).
\begin{itemize}
  \item Non c'è coupling con il database del microservizio di origine.
  \item C'è solo un legame con gli eventi emessi dal servizio, che sono fatti apposta per essere esposti.
  \item Il flusso di dati basato sugli eventi arriva al sistema di reporting più in fretta di un data pump eseguito a
    intervalli fissi.
  \item L'event data pump non deve per forza essere gestito dal team del servizio.
  \item Svantaggio: tutte le informazioni necessarie devono essere diffuse come eventi. Con grandi volumi di dati la
    soluzione può non scalare bene.
\end{itemize}

\begin{figure}[ht]
\centering
\begin{tikzpicture}
\node[sv, align=center] (cs) at (0,0) {Customer\\service};
\node[circle, draw, fill=arancio!60, minimum size=0.42cm] (ev) at (2.0,0) {};
\node[font=\footnotesize, align=center, anchor=north] at (2.0,-0.3) {eventi di\\cambio di stato};
\node[draw, minimum width=0.7cm, minimum height=1.9cm] (q) at (3.6,0) {};
\foreach \y in {-0.6,0,0.6} \node[circle, draw, fill=arancio!60, minimum size=0.42cm] at (3.6,\y) {};
\node[sv, align=center, minimum width=2.4cm] (mp) at (6.2,0) {Customer-reporting\\mapper};
\node[dbs] (rd) at (9.4,0) {Central\\reporting\\database};
\draw (cs) -- (ev);
\draw[->] (ev) -- (q);
\draw[->] (q) -- (mp);
\draw[->] (mp) -- (rd);
\end{tikzpicture}
\caption{Event data pump. Il servizio Customer emette un evento a ogni cambio di stato di un cliente. Gli eventi passano da
una coda; il mapper, che è il subscriber, li legge e aggiorna il database centrale di reporting.}
\label{fig:split-event}
\end{figure}

\begin{table}[ht]
\centering
\small
\caption{Confronto delle tre soluzioni per il reporting.\aggiunta}
\label{tab:split-reporting}
\begin{tabularx}{\linewidth}{>{\raggedright\arraybackslash}p{2.9cm}>{\raggedright\arraybackslash}p{2.6cm}>{\raggedright\arraybackslash}X>{\raggedright\arraybackslash}p{3.2cm}}
\toprule
Soluzione & Chi muove i dati & Legame con il servizio & Limite \\
\midrule
Chiamate ai servizi & il reporting (pull) & solo l'API pubblica & grandi volumi; API non adatte \\
Data pump & il data pump (push) & schema interno del database; stesso team e versione & coupling con il database \\
Event data pump & il subscriber degli eventi & solo gli eventi pubblici & tutto deve diventare evento; grandi volumi \\
\bottomrule
\end{tabularx}
\end{table}

\section{Considerazioni finali}
\label{sec:split-conclusioni}

Scomporre un monolite può essere faticoso.
\begin{itemize}
  \item Pianificare cambiamenti \textbf{piccoli e incrementali} rende più facile capire l'impatto di ogni cambiamento.
    Rende anche più facile contenere il costo degli errori, che sono inevitabili.
  \item Far crescere un servizio finché va diviso va del tutto bene. Bisogna però capire \emph{quando} va diviso, e
    dividerlo, prima che la divisione diventi troppo costosa.
\end{itemize}

Riferimento: Sam Newman, \emph{Building Microservices}, O'Reilly, 2021, capitolo 5.

\begin{puntichiave}
\begin{itemize}
  \item Crescendo, il monolite perde cohesion e loose coupling. Si scompone quando diventa un problema.
  \item Un seam è una porzione di codice trattabile in isolamento. I seam diventano confini di servizio: si trovano con i
    bounded context e si materializzano come package, in modo incrementale.
  \item Si inizia dal seam con più beneficio: pace of change, team structure, security, technology, tangled dependencies
    (il seam meno dipeso).
  \item Il database è la dipendenza più aggrovigliata. Refactoring tipici:
    \begin{itemize}
      \item foreign key tra contesti $\to$ API del servizio proprietario (vincoli gestiti nei servizi);
      \item dati statici condivisi $\to$ copia, configurazione (preferita) o servizio;
      \item dati modificabili condivisi $\to$ nuovo servizio per il concetto;
      \item tabella con due concetti $\to$ due tabelle.
    \end{itemize}
  \item Separare prima gli schemi, con il codice ancora insieme, rende più facile tornare indietro.
  \item Transazioni su più schemi: transazioni compensative, two-phase commit (con le sue vulnerabilità) o riprovare più
    tardi con eventual consistency.
  \item Reporting su più servizi: chiamate alle API, data pump (gestito dal team del servizio), event data pump.
\end{itemize}
\end{puntichiave}

\begin{domande}
\begin{enumerate}
  \item Perché un monolite che cresce perde cohesion e loose coupling? Quando conviene scomporlo?
  \item Che cos'è un seam? Come aiutano i bounded context e i package a trovarlo?
  \item Quali sono i criteri per scegliere il primo seam da estrarre? Fare un esempio per ciascuno.
  \item Perché il database è la dipendenza più difficile da spezzare? Quali strumenti aiutano ad analizzarlo?
  \item Come si elimina una foreign key tra le tabelle di due contesti? Che cosa si perde?
  \item Quali soluzioni ci sono per i dati statici condivisi? E per i dati modificabili condivisi?
  \item Conviene separare prima gli schemi o prima i servizi? Perché?
  \item Che cosa succede alle transazioni quando si separano gli schemi? Confrontare le tre soluzioni.
  \item Descrivere il two-phase commit e le sue vulnerabilità.
  \item Come si fa reporting quando il database è diviso? Confrontare data pump ed event data pump.
\end{enumerate}
\end{domande}
